\documentclass[11pt]{article}

\usepackage[margin=0.8in]{geometry}
\usepackage{amsmath,amssymb}
\usepackage{booktabs}
\usepackage{enumitem}
\usepackage[hidelinks]{hyperref}

\setlist[itemize]{leftmargin=1.4em}
\setlist[enumerate]{leftmargin=1.6em}

\title{SYSC 5108 Exam Study\\Assignment 4, Question 5}
\author{Jesse Levine}
\date{}

\begin{document}
\maketitle

\section*{Question Summary}

We have a two-state, two-action reinforcement-learning problem:

\begin{center}
\begin{tabular}{ccccc}
\toprule
State & Action & Next state & Probability & Reward \\
\midrule
0 & 0 & 0 & 1 & -1 \\
0 & 1 & 1 & 1 & 1 \\
1 & 0 & 1 & 1 & 1 \\
1 & 1 & 0 & 1 & -1 \\
\bottomrule
\end{tabular}
\end{center}

Assume:
\[
    \gamma=0.1,
    \qquad
    \pi(1\mid 0)=\pi(0\mid 1)=0.6.
\]

The question asks for:
\begin{enumerate}[label=(\alph*)]
    \item \(\varepsilon\) for the \(\varepsilon\)-greedy policy,
    \item the state-value and action-value functions under the given policy,
    \item the optimal policy and optimal values,
    \item Monte Carlo estimates from two generated trajectories,
    \item off-policy TD control (Q-learning),
    \item on-policy TD control (SARSA).
\end{enumerate}

\section*{Course and Textbook Cross-References}

\begin{itemize}
    \item \textbf{Chapter 6 slides, Value Function and Solving Bellman Equation:}
    use Bellman equations for \(v_\pi(s)\) and \(q_\pi(s,a)\).
    \item \textbf{Chapter 6 slides, \(\varepsilon\)-greedy policy:} the greedy
    action gets probability \(1-\varepsilon+\varepsilon/|\mathcal{A}|\).
    \item \textbf{Chapter 6 slides, Monte Carlo:} update values from returns
    \(G_t\).
    \item \textbf{Chapter 6 slides, Q-learning:}
    \[
        q(s,a)\leftarrow q(s,a)+\alpha\left[r+\gamma\max_{a'}q(s',a')-q(s,a)\right].
    \]
    \item \textbf{Chapter 6 slides, SARSA:}
    \[
        q(s,a)\leftarrow q(s,a)+\alpha\left[r+\gamma q(s',a')-q(s,a)\right].
    \]
\end{itemize}

\section*{Part (a): Find \(\varepsilon\)}

There are two actions, so an \(\varepsilon\)-greedy policy gives the greedy action
probability
\[
    1-\varepsilon+\frac{\varepsilon}{2}=1-\frac{\varepsilon}{2}.
\]
The problem states
\[
    \pi(1\mid 0)=\pi(0\mid 1)=0.6,
\]
and these are the greedy-action probabilities in states \(0\) and \(1\),
respectively. Therefore
\[
    1-\frac{\varepsilon}{2}=0.6
    \quad\Rightarrow\quad
    \frac{\varepsilon}{2}=0.4
    \quad\Rightarrow\quad
    \varepsilon=0.8.
\]

\[
    \boxed{\varepsilon=0.8}
\]

\section*{Part (b): Find \(V_\pi\) and \(Q_\pi\)}

Under the given policy:
\[
    \pi(0\mid 0)=0.4,\qquad \pi(1\mid 0)=0.6,
\]
\[
    \pi(0\mid 1)=0.6,\qquad \pi(1\mid 1)=0.4.
\]

The Bellman equations are
\[
    V(0)=0.4[-1+0.1V(0)] + 0.6[1+0.1V(1)],
\]
\[
    V(1)=0.6[1+0.1V(1)] + 0.4[-1+0.1V(0)].
\]
By symmetry, let
\[
    V(0)=V(1)=v.
\]
Then
\[
\begin{aligned}
    v
    &=0.4(-1+0.1v)+0.6(1+0.1v) \\
    &=-0.4+0.04v+0.6+0.06v \\
    &=0.2+0.1v.
\end{aligned}
\]
So
\[
    0.9v=0.2
    \quad\Rightarrow\quad
    v=\frac{2}{9}\approx 0.2222.
\]

Therefore
\[
    \boxed{
    V(0)=V(1)=\frac{2}{9}\approx 0.2222
    }.
\]

Now compute the action values:
\[
    Q(0,0)=-1+0.1V(0)=-1+\frac{1}{45}=-\frac{44}{45}\approx -0.9778,
\]
\[
    Q(0,1)=1+0.1V(1)=1+\frac{1}{45}=\frac{46}{45}\approx 1.0222,
\]
\[
    Q(1,0)=1+0.1V(1)=\frac{46}{45}\approx 1.0222,
\]
\[
    Q(1,1)=-1+0.1V(0)=-\frac{44}{45}\approx -0.9778.
\]

\[
    \boxed{
    \begin{aligned}
    Q(0,0)&=-0.9778, & Q(0,1)&=1.0222,\\
    Q(1,0)&=1.0222,  & Q(1,1)&=-0.9778.
    \end{aligned}}
\]

\section*{Part (c): Find the Optimal Policy}

The optimal Bellman equations are
\[
    V^*(0)=\max\{-1+0.1V^*(0),\;1+0.1V^*(1)\},
\]
\[
    V^*(1)=\max\{1+0.1V^*(1),\;-1+0.1V^*(0)\}.
\]
Clearly the better action in state \(0\) is action \(1\), and the better action in
state \(1\) is action \(0\). Thus
\[
    \boxed{\pi^*(0)=1,\qquad \pi^*(1)=0}.
\]

Under this policy,
\[
    V^*(1)=1+0.1V^*(1)
    \quad\Rightarrow\quad
    0.9V^*(1)=1
    \quad\Rightarrow\quad
    V^*(1)=\frac{10}{9}.
\]
Then
\[
    V^*(0)=1+0.1V^*(1)=1+\frac{1}{9}=\frac{10}{9}.
\]

\[
    \boxed{
    V^*(0)=V^*(1)=\frac{10}{9}\approx 1.1111
    }.
\]

\section*{Part (d): Monte Carlo Estimates}

Use the random numbers in order:
\[
0.94,\;0.27,\;0.30,\;0.18,\;0.75,\;0.35,\ldots
\]

Action-selection convention under the given \(\varepsilon\)-greedy policy:
\begin{itemize}
    \item in state \(0\), choose action \(1\) if \(r\le 0.6\), else action \(0\),
    \item in state \(1\), choose action \(0\) if \(r\le 0.6\), else action \(1\).
\end{itemize}

\subsection*{Trajectory starting from \(s_0=0\)}

\[
    0 \xrightarrow[a=0]{-1} 0 \xrightarrow[a=1]{1} 1 \xrightarrow[a=0]{1} 1.
\]
This uses random numbers \(0.94, 0.27, 0.30\).

\subsection*{Trajectory starting from \(s_0=1\)}

\[
    1 \xrightarrow[a=0]{1} 1 \xrightarrow[a=1]{-1} 0 \xrightarrow[a=1]{1} 1.
\]
This uses random numbers \(0.18, 0.75, 0.35\).

\subsection*{Returns}

For the first trajectory:
\[
    G_0^{(1)}=-1+0.1(1)+0.1^2(1)=-0.89,
\]
\[
    G_1^{(1)}=1+0.1(1)=1.1,
\]
\[
    G_2^{(1)}=1.
\]

For the second trajectory:
\[
    G_0^{(2)}=1+0.1(-1)+0.1^2(1)=0.91,
\]
\[
    G_1^{(2)}=-1+0.1(1)=-0.9,
\]
\[
    G_2^{(2)}=1.
\]

\subsection*{State-value estimates}

Using every-visit incremental Monte Carlo with \(\alpha=0.6\) and initial values
equal to \(0\),
\[
    \boxed{
    V(0)\approx 0.77856,\qquad V(1)\approx -0.22560
    }.
\]

\subsection*{Action-value estimates}

Again using every-visit incremental Monte Carlo with \(\alpha=0.6\):
\[
    \boxed{
    Q(0,0)=-0.534,\quad
    Q(0,1)=0.864,\quad
    Q(1,0)=0.786,\quad
    Q(1,1)=-0.540
    }.
\]

\section*{Part (e): Off-Policy TD Control (Q-learning)}

Use
\[
    Q(s,a)\leftarrow Q(s,a)+\alpha\left[r+\gamma\max_{a'}Q(s',a')-Q(s,a)\right],
\]
with \(\alpha=0.6\), starting from
\[
    Q(s,a)=0
    \quad\text{for all }(s,a).
\]
For the final step of each trajectory, treat the endpoint at \(t=3\) as terminal, so
the target is just the immediate reward.

After processing the two trajectories from part (d), the final Q-values are:
\[
    \boxed{
    Q(0,0)=-0.600,\quad
    Q(0,1)=0.840,\quad
    Q(1,0)=0.876,\quad
    Q(1,1)=-0.564
    }.
\]

\section*{Part (f): On-Policy TD Control (SARSA)}

Use
\[
    Q(s,a)\leftarrow Q(s,a)+\alpha\left[r+\gamma Q(s',a')-Q(s,a)\right].
\]
Start with
\[
    Q(s,a)=0
    \quad\text{for all }(s,a),
\]
and initial policy
\[
    \pi(1\mid 0)=\pi(0\mid 1)=0.5.
\]

Using the same random sequence from the beginning,
\[
0.94,\;0.27,\;0.30,\;0.18,\;0.75,\ldots
\]
the trajectory is
\[
    0 \xrightarrow[a=0]{-1} 0
    \xrightarrow[a=1]{1} 1
    \xrightarrow[a=0]{1} 1
    \xrightarrow[a=0]{1} 1
    \xrightarrow[a=1]{-1} 0.
\]

Applying SARSA step by step gives:
\[
    \boxed{
    Q(0,0)=-0.600,\quad
    Q(0,1)=0.600,\quad
    Q(1,0)=0.840,\quad
    Q(1,1)=-0.600
    }.
\]

\section*{Final Answer}

\[
    \boxed{\varepsilon=0.8}
\]

\[
    \boxed{
    V(0)=V(1)=\frac{2}{9}\approx 0.2222
    }
\]

\[
    \boxed{
    \begin{aligned}
    Q(0,0)&=-0.9778, & Q(0,1)&=1.0222,\\
    Q(1,0)&=1.0222,  & Q(1,1)&=-0.9778
    \end{aligned}}
\]

\[
    \boxed{\pi^*(0)=1,\qquad \pi^*(1)=0}
\]

\[
    \boxed{
    V^*(0)=V^*(1)=\frac{10}{9}\approx 1.1111
    }
\]

\[
    \boxed{
    \text{MC: }
    V(0)\approx 0.77856,\;
    V(1)\approx -0.22560
    }
\]

\[
    \boxed{
    \text{MC }Q:
    Q(0,0)=-0.534,\;
    Q(0,1)=0.864,\;
    Q(1,0)=0.786,\;
    Q(1,1)=-0.540
    }
\]

\[
    \boxed{
    \text{Q-learning: }
    Q(0,0)=-0.600,\;
    Q(0,1)=0.840,\;
    Q(1,0)=0.876,\;
    Q(1,1)=-0.564
    }
\]

\[
    \boxed{
    \text{SARSA: }
    Q(0,0)=-0.600,\;
    Q(0,1)=0.600,\;
    Q(1,0)=0.840,\;
    Q(1,1)=-0.600
    }
\]

\section*{Exam Notes}

\begin{itemize}
    \item In this problem, the transitions are deterministic. The only randomness
    is from the behavior policy.
    \item For \(\varepsilon\)-greedy with two actions, the greedy action gets
    probability \(1-\varepsilon/2\), not \(1-\varepsilon\).
    \item Q-learning uses the greedy target
    \(\max_{a'}Q(s',a')\), while SARSA uses the actually selected next action
    \(a'\).
    \item Be explicit about how the random numbers are mapped to actions. That is
    where most bookkeeping mistakes happen.
\end{itemize}

\end{document}
